\section{Introduction and Literature Survey}
\label{sec:intro}

\subsection{Background and motivation}
\label{ssec:motivation}

Underwater acoustic (UWA) communications carry much of the non-tethered data
traffic between autonomous underwater vehicles, ocean-bottom sensors, and surface
platforms. Their physical layer is shaped by two channel pathologies that, taken
together, are harder to handle than almost any terrestrial radio-frequency link.
The multipath between the sea surface, the water column, and the sea bed produces
delay spreads of several to tens of milliseconds, which at the 4--12~kHz bandwidth
budgets typical of UWA modems translates to hundreds of resolvable delay
bins~\cite{stojanovic2009uwa,preisig2007uwa}. The propagation speed of $c \approx
1500$~m/s, combined with platform motion of even a few~m/s, yields a relative
Doppler of order $v / c \sim 10^{-3}$~--- three to four orders of magnitude larger
than in a typical RF link --- and produces both strong inter-carrier interference
and fractional Doppler offsets relative to the underlying multicarrier
grid~\cite{berger2010uwasparse,huang2010uwamimoofdm}. A conventional OFDM receiver
cannot absorb either pathology cleanly: a cyclic prefix long enough for the full
delay spread costs a prohibitive fraction of spectral efficiency, while a residual
fractional Doppler of even a few tenths of a sub-carrier corrupts sub-carrier
orthogonality in a way that a single-tap frequency-domain equaliser cannot repair.

Affine Frequency Division Multiplexing (AFDM)~\cite{bemani2023afdm} was designed
for precisely such doubly dispersive channels. By placing information on a chirp
basis defined through the discrete affine Fourier transform (DAFT), and choosing
the two chirp coefficients $(c_1, c_2)$ to align with the channel's
delay--Doppler geometry, AFDM folds every physical delay--Doppler pair into a
single sparse structured effective channel on the chirp-domain received signal,
rather than into the full two-dimensional delay--Doppler grid that Orthogonal
Time Frequency Space (OTFS) modulation~\cite{hadani2017otfs,raviteja2018otfs,%
wei2021otfspromising} would produce. The resulting sparsity is what makes
sparse-recovery channel estimation with orthogonal matching pursuit
(OMP)~\cite{tropp2007signal} tractable~\cite{benzine2023sparseafdm,%
rou2024otfstoafdm}, and is the starting point of every receiver considered in
this report.

Pilot overhead nonetheless remains a first-order constraint on spectral
efficiency in UWA contexts where the bandwidth budget is only a few kilohertz.
To preserve the data-carrying capacity that an AFDM frame would otherwise lose
to dedicated pilots, this work adopts a superimposed-pilot (SP)
transmission~\cite{orozco2011superimposed,tugnait2006superimposed} in which
every chirp sub-carrier simultaneously carries data and pilot energy under a
fixed power split $\rhop$. SP doubles the available data capacity relative to a
comparable embedded-pilot (EP) frame, but it also introduces a new estimation
problem: the pilot observation is now systematically contaminated by
unknown-data interference, and at mid-to-high SNR this interference does not
average out. In a UWA-style $\rhop = 0.5$, $\Ppath = 8$-path simulation
representative of the regime considered here, the resulting error floor for a
one-shot OMP + MP-BP receiver sits at $\ber \approx 5 \times 10^{-2}$ from
$\snr \geq 15$~dB, with an OMP channel-estimation NMSE frozen near $-11.3$~dB
irrespective of the noise variance. The presence of this floor --- not the SNR
--- defines the engineering regime that this work has to address.

The natural response to a structured estimation residual is to iterate between
channel estimation and symbol detection~\cite{wymeersch2007iterativereceiver}:
each detection pass produces tentative data estimates that can be subtracted
from the pilot observation, sharpening the next channel estimate. Two naive
realisations of this idea are expensive, however. Re-running OMP at every outer
iteration multiplies the per-frame cost by a 5$\times$ constant factor at five
outer iterations, and replacing OMP with a dense LMMSE channel estimator
inflates the per-frame complexity by two orders of magnitude at $\Nsub = 1024$.
This work therefore investigates a third option: estimate the channel support
\emph{once} with OMP, lock that support, and iterate only the cheap downstream
blocks --- a damped data-aided least-squares refit of the path gains and sparse
message-passing belief propagation (MP-BP) on the chirp-domain factor graph.
The central engineering question is whether such a single-OMP receiver,
optionally augmented with a small learned module that repairs the support
residual that fractional-Doppler leakage leaves behind, can close the SP error
floor at a complexity comparable to a single OMP pass.


\subsection{Critical literature survey}
\label{ssec:lit}

The proposed receiver sits at the intersection of four research streams that
have, until now, evolved largely in isolation from each other: AFDM as a
chirp-domain waveform engineered for doubly dispersive channels; superimposed
pilot transmission as a spectral-efficiency mechanism whose joint
estimation--detection challenge predates AFDM by nearly two decades;
sparse-recovery and message-passing receiver primitives developed in the
OFDM/OTFS lineage and now inherited by AFDM; and the deep-unfolding tradition
that motivates the embedding of small learned modules within otherwise
model-based receivers. This section traces each thread to the point at which
it constrains a specific design choice in Chapter~\ref{sec:method}, and
concludes by identifying the narrow gap that the present work occupies.

\paragraph{AFDM as a doubly-dispersive waveform and its channel-estimation
primitives.}
AFDM was introduced by Bemani \emph{et al.}~\cite{bemani2023afdm} as a
chirp-based multicarrier waveform whose two tunable DAFT parameters $(c_1,
c_2)$ can be aligned with the delay--Doppler geometry of the underlying
channel, yielding a sparse one-dimensional input--output relation that
recovers the diversity advantages of the two-dimensional OTFS
modulation~\cite{hadani2017otfs,raviteja2018otfs,wei2021otfspromising} on
doubly dispersive links. Subsequent investigations established the
receiver-side primitives that make this representation operational. Yin and
Tang~\cite{yin2022fractionaldopplerafdm} developed the pilot-aided AFDM
channel estimator for the integer delay--Doppler grid, formalising the
threshold-based recovery rule that later sparse and Bayesian extensions would
refine. Benzine \emph{et al.}~\cite{benzine2023sparseafdm} reformulated AFDM
channel estimation as a sparse-recovery problem under embedded-pilot (EP)
guard intervals, importing OMP~\cite{tropp2007signal} into the AFDM
literature and demonstrating that the sparsity of the chirp-domain effective
channel makes greedy recovery a natural fit. Yin \emph{et
al.}~\cite{yin2024drceafdm} extended this line to MIMO-AFDM by deriving a
diagonal-reconstruction estimator that explicitly addresses the inter-Doppler
interference produced by fractional Doppler residuals, while the comparative
survey of Rou \emph{et al.}~\cite{rou2024otfstoafdm} situated AFDM alongside
OTFS and OCDM within an integrated sensing-and-communication context. Two
practitioner-oriented overviews~\cite{yin2025afdmoverview,zhou2025afdmicsc}
now collect these threads. Across this body of work, the EP frame keeps the
pilot observation pilot-only and yields analytically clean sparse-recovery
guarantees, but it does so at a spectral-efficiency cost that is hard to
absorb at the few-kilohertz bandwidth budgets typical of UWA modems, and the
analytical tractability disappears once the EP assumption is dropped.

\paragraph{Superimposed pilots and the joint estimation--detection challenge.}
The superimposed-pilot principle, in which every symbol simultaneously
carries data and pilot energy under a fixed power split, was developed for
OFDM more than a decade before AFDM. Cui and
Tellambura~\cite{cui2007superimposed} derived ML and MMSE iterative
estimators for SP-OFDM and proved that the equispaced equipowered placement
minimises the Cram\'er--Rao bound, while Tugnait and
Meng~\cite{tugnait2006superimposed} characterised SP-training
power-allocation tradeoffs and frame-synchronisation behaviour. Orozco-Lugo
\emph{et al.}~\cite{orozco2011superimposed} established the optimality of
the equipowered training pattern under the spectral-efficiency constraint. Liang
\emph{et al.}~\cite{liang2011iterativesp} subsequently demonstrated that
superimposed pilots can match preamble-based channel estimation in OFDM-WLAN
at the cost of an iterative receiver, isolating the regime in which the
iteration overhead is justified by the bandwidth saving. The structural
observation that survives the transfer from OFDM to delay--Doppler waveforms
is that the pilot--data interference does not decay with the noise variance
and therefore produces a high-SNR error floor that single-pass estimation
cannot remove. In delay--Doppler waveforms specifically, Mishra \emph{et
al.}~\cite{mishra2022otfssp} extended the SP framework to OTFS, and Neelam
and Sahu~\cite{neelam2023otsmsp} carried it to OTSM, in both cases confirming
that an iterative loop alternating pilot-data interference cancellation and
channel re-estimation is unavoidable once the EP assumption is abandoned.

The most directly relevant SP-AFDM line is the recent cluster initiated by
Zheng \emph{et al.}~\cite{qing2024spafdm}, which adapts the OFDM SP machinery
to the chirp domain and introduces an iterative SP-AFDM channel estimator
and signal detector with optimised pilot placement. Cao \emph{et
al.}~\cite{cao2024spomp} subsequently combined the SP frame with an OMP
front-end, leveraging the chirp-domain sparsity to lower the per-iteration
estimation cost. Both works establish that, under SP-AFDM, an iterative loop
is structurally required, but they re-execute the sparse-recovery step on
every outer iteration and contain no learning component, leaving open the
question of whether the support recovered by the first OMP pass already
captures the bulk of the channel energy and therefore makes per-iteration
re-OMP redundant. Adjacent to these SP-focused works, Zhou \emph{et
al.}~\cite{zhao2024gifreeafdm} reach a structurally similar conclusion via a
guard-interval-free pilot-aided AFDM design with iterative interference
cancellation; their formulation, however, departs from SP in that it
preserves a separation between pilot and data symbols and addresses only
data-to-pilot and pilot-to-data interference rather than the multiplicative
pilot--data coupling that defines the SP frame. Huang \emph{et
al.}~\cite{afdm2025jointcesd} propose AFDM channel-estimation and
symbol-detection receivers including a DNN-based design and a separate
GI-free iterative design, demonstrating that learning-augmented AFDM
receivers are tractable in principle, but their learned element is placed at
the detector rather than at the channel-support-selection step, and the work
targets neither an SP frame nor a UWA-style channel model. Across this cluster, the iterative joint-receiver
philosophy is now well established within AFDM, yet no published
contribution simultaneously preserves the SP framework, retains the
sparse-recovery skeleton without re-running OMP, and places a small learned
module at the OMP support-refinement step.

\paragraph{Sparse channel estimation in UWA contexts and physical-layer
refinements.}
The choice of OMP as the sparse-recovery primitive in this work inherits a
long line of evidence that UWA channels are well approximated by a sparse
delay--Doppler representation. Stojanovic and
Preisig~\cite{stojanovic2009uwa} characterised the underwater acoustic
channel as severely doubly dispersive with a sparse impulse response in
which each physical path acts as a time-varying low-pass filter, while
Preisig~\cite{preisig2007uwa} formalised the propagation considerations
governing UWA system design. Berger \emph{et al.}~\cite{berger2010uwasparse}
then established compressive sensing as the natural channel-estimation
primitive for multicarrier UWA links, and Huang \emph{et
al.}~\cite{huang2010uwamimoofdm} extended the sparse-recovery toolkit to
MIMO-OFDM-UWA with iterative channel estimation and decoding. Within the
OFDM-UWA literature this lineage diversified into greedy variants such as
backtracking OMP~\cite{junejo2018labomp} and Bayesian variants based on
sparse Bayesian learning, all of which testify to the practical robustness
of the sparse-recovery approach but operate within OFDM rather than within a
chirp-domain waveform.

The closest physical-layer refinements to the AFDM-UWA setting considered
here form a cluster of three threads. Ranasinghe \emph{et
al.}~\cite{ma2024oversampledafdm} show that observing the chirp-domain signal
at twice the Nyquist rate yields a tangible channel-estimation improvement
that is independent of any learned component; this directly motivates the
$\fs = 4B$ operating point adopted in the present report, although their
result is established in a non-SP, non-learning configuration and does not
address either pilot-data interference or fractional-Doppler support
residuals. Cao \emph{et al.}~\cite{zhou2024msmlafdm} extend AFDM-OMP to
multi-scale multi-lag (MSML) UWA channels in which each physical path
carries an independent time-scale factor, establishing that the AFDM
sparse-recovery primitive transfers to a more physically faithful UWA
channel model; their work, however, is non-SP and contains no learning
component, and therefore engages with neither the SP-induced estimation
floor nor the Dirichlet-leakage support residual that motivates the present
design. The fractional-Doppler problem itself has been treated directly by
Yin \emph{et al.}~\cite{yin2024drceafdm} via a diagonal-reconstruction
estimator targeting the very Dirichlet leakage that single-pass OMP fails to
capture, but the formulation is non-SP and the leakage repair is not
embedded inside an iterative joint receiver. Together these three threads
establish the individual physical-layer ingredients --- oversampling, MSML
treatment, Dirichlet-leakage repair --- yet none is integrated with the
SP-induced estimation-limited floor that this report is designed to close.

\paragraph{Sparse message-passing detection.}
On the detection side, the proposed inner LS/MP-BP loop inherits from the
sparse delay--Doppler message-passing literature pioneered by Raviteja
\emph{et al.}~\cite{raviteja2018otfs}, in which a low-complexity
factor-graph detector exploits the natural sparsity of the OTFS channel
matrix and explicitly handles the inter-Doppler interference produced by
fractional Doppler shifts. Subsequent OTFS detector designs --- including
unitary-AMP variants~\cite{yuan2022uampotfs}, expectation-propagation
formulations~\cite{pu2024epotfs}, and graph-neural-network
refinements~\cite{zhuang2024amgnnotfs} --- have refined the
complexity--accuracy tradeoff but preserve the sparse-graph philosophy. The
present work transports this MP-BP machinery into the chirp-domain factor
graph induced by the Dirichlet-expanded AFDM virtual-tap representation,
and uses it as a fixed inner-loop detector throughout.

\paragraph{Deep unfolding for model-based receivers.}
At the architectural level, the proposed learned head sits within the
deep-unfolding tradition surveyed by Monga \emph{et
al.}~\cite{monga2021unrolling} and Balatsoukas-Stimming and
Studer~\cite{balatsoukasstimming2019deepunfolding}, and rooted in the
seminal Learned ISTA architecture of Gregor and
LeCun~\cite{gregor2010lista}, in which iterations of a sparse-recovery
algorithm are unrolled into trainable layers; the AMP-Inspired Deep
Networks of Borgerding \emph{et al.}~\cite{borgerding2017amp} extend this
template to approximate-message-passing iterations and demonstrate that
unfolded sparse-recovery networks recover the linear-inverse-problem
performance of their model-based counterparts at a fraction of the
iteration count. The model-driven survey of He
\emph{et al.}~\cite{he2019modeldrivensurvey} situates this philosophy
within physical-layer communications, arguing that combining domain-derived
structure with learned components reduces both the training-data demand and
the inference complexity relative to fully data-driven receivers. The bulk
of the resulting literature targets MIMO detection --- exemplified by the
model-driven detectors of He \emph{et
al.}~\cite{he2020modeldrivenmimo}, the soft interference
cancellation network DeepSIC~\cite{shlezinger2021deepsic}, the
adaptive-detection scheme MMNet~\cite{khani2020mmnet}, and the
attention-augmented receiver of Xu \emph{et
al.}~\cite{xu2024detecttolearn}. Within AFDM specifically, learned
components have so far appeared at the detector~\cite{afdm2025jointcesd}
rather than at the channel-support-selection step, and within the broader
sparse-recovery literature the unrolling of greedy-OMP iterations has been
studied only recently~\cite{mohammadtaheri2025deepgreedyunfolding}, where
the non-differentiable argsort operator is approached via softsort
relaxations. The OMP support-refinement slot inside an SP-AFDM joint
receiver, however, remains untouched by learning.

\paragraph{Synthesis and identified gap.}
Across the four threads above, the design space for an SP-AFDM receiver in
a UWA-style sparse fractional-Doppler regime leaves a narrow opening. The
receiver should preserve the spectral-efficiency advantage of SP and the
sparse-recovery prior that makes AFDM attractive in UWA, while avoiding
both the per-iteration re-OMP overhead that the existing SP-AFDM iterative
receivers incur and the cubic per-frame cost of a dense LMMSE estimator.
It should additionally address the residual fractional-Doppler support
errors that classical OMP leaves behind, ideally without enlarging the
dictionary or altering the asymptotic complexity class. To the author's
knowledge, no published AFDM receiver simultaneously runs OMP only once,
iterates only over a damped data-aided least-squares path-gain refit and
sparse MP-BP detection, and embeds a small learned module dedicated to OMP
support refinement under the UWA-style sparse SP regime. The receiver
developed in this report is positioned precisely at this intersection.


\subsection{Research question and objectives}
\label{ssec:rq}

The central research question of the report is:
\begin{quote}
\emph{Can an SP-AFDM receiver close the structured pilot-data interference
floor at mid-to-high SNR while remaining complexity-preserving relative to both
re-OMP-per-iteration and dense-LMMSE alternatives, and can a minimal learned
module further reduce the residual support-selection error without leaving the
sparse-recovery complexity class?}
\end{quote}

This question is decomposed into three operational objectives. The first is to
characterise the SP-induced estimation floor and demonstrate quantitatively
that a one-shot OMP + MP-BP receiver is estimation-limited rather than
noise-limited. The second is to design and justify an OMP-once + iterative
LS/MP-BP receiver against both its re-OMP-per-iteration overhead variant and
its dense-LMMSE alternative, and to verify the design choice with a direct
support-recall measurement. The third is to design a learned
support-refinement module that strictly preserves the classical
sparse-recovery complexity class, and to quantify its BER contribution under
a paired-frame statistical protocol.


\subsection{Contributions}
\label{ssec:contrib}

The report makes three contributions, each summarised below.

\begin{enumerate}[leftmargin=1.3em,itemsep=2pt]
\item \textbf{A complexity-preserving joint OMP--MP receiver for AFDM-SP under a
UWA-style fractional-Doppler regime.} The proposed receiver runs OMP only
once, locks the recovered support, and iterates only the cheap downstream
blocks: a damped data-aided least-squares refit of the path gains and sparse
MP-BP detection on the chirp-domain factor graph. A direct support-recall
probe shows that the first OMP pass already recovers a mean $\approx 95\%$ of
the effective Dirichlet-expanded virtual-tap support, with the residual
$\approx 5\%$ stable across SNR and therefore structurally rather than
noise-driven; re-running OMP across outer iterations is consequently a wasted
constant-factor overhead. The proposed receiver recovers two to three orders
of magnitude in BER over the one-shot OMP + MP-BP baseline at $\snr = 25$~dB,
while staying within the sparse $\mathcal{O}(\Ppath \Nsub^2)$ arithmetic
class.

\item \textbf{A 1\,041-parameter learned support-refinement head deployed as a
deterministic top-$\Kadd$ ranker.} A small deep-unfolding head, inserted
between the single OMP pass and the iterative LS/MP-BP loop, scores
off-support candidates and appends a fixed number of top-ranked atoms to the
active support. Pooled across three independent evaluation-seed blocks, the
head delivers a consistent $\approx 2$~dB paired BER separation from the
classical joint receiver at $\snr \in \{20, 25, 30\}$~dB, while keeping the
receiver in the same $\mathcal{O}(\Ppath \Nsub^2)$ arithmetic class as the
classical baseline. The analytical complex-MAC overhead at $\Nsub = 1024$ is
$\approx 8\%$ over the classical core; on a CPU/GPU reference,
the measured wall-clock latency is at parity with the classical-Torch baseline
at $\Nsub = 1024$ and within $+3.9\%$ at $\Nsub = 2048$, so the analytical
overhead does not reverse in practice.

\item \textbf{A mechanism analysis grounding the learned gain.} A
top-$\Ppath$ physical-support-oracle diagnostic provides a valid ceiling on
the BER reduction available from any support-refinement strategy operating on
the proposed dictionary. Against this ceiling, the learned head closes
$61\%/81\%/98\%$ of the achievable gap at $\snr = 15/20/25$~dB. Combined with
the support-recall probe of Contribution~1, this attributes the learned
head's BER contribution specifically to the recovery of Dirichlet side-lobe
atoms that single-pass OMP discards under fractional-Doppler leakage, rather
than to a generic noise-averaging effect.
\end{enumerate}


\subsection{Report structure}
\label{ssec:structure}

Chapter~\ref{sec:method} develops the AFDM-SP system model, the OMP
dictionary derivation, the proposed OMP-once + iterative LS/MP-BP recipe, the
learned support-refinement head with its training and inference setup, and
the per-frame complexity analysis. It also documents the paired-frame
evaluation protocol and the choice of baselines.
Chapter~\ref{sec:results} reports four experiments: the main BER waterfall,
the outer-iteration convergence probe, the design-rationale support-recovery
probes, and the complexity comparison.
Chapter~\ref{sec:discussion} interprets the results mechanistically, situates
them against the prior AFDM literature, and discusses methodological and
scope limitations.
Chapter~\ref{sec:conclusion} summarises the conclusions and outlines concrete
follow-up work.
